\documentclass{situationist-preprint}

\title{Substrata and Individuation}
\msid{MET-23-b}
\msfamily{MET-23}
\msdiscipline{Metaphysics}
\mstype{Consequence}
\mskeywords{individuation; bare particulars; substrata; haecceity; numerical distinctness; weak discernibility}

\begin{document}
\maketitle

\begin{abstract}
The principal motive for bare substrata has been individuation: something must make two qualitatively alike objects two. MET-23-a concludes that there are no bare particulars. We derive what follows for individuation. Numerical distinctness is not conferred by a constituent; it is a fact about the objects themselves. Qualities do not individuate in the absolute or relative sense, by MET-28, and real definitions do not individuate either. What remains is primitive distinctness, which irreflexive relations may witness, as they weakly discern Black's spheres, and which haecceitistic properties presuppose rather than supply. Diachronic individuation is likewise not by substrata but by the persistence conditions of kinds.
\end{abstract}

\section{What does not individuate}

Call a theory of individuation an answer to the question in virtue of what two objects are two. Theorem~3.3 of MET-23-a, that there are no bare particulars, rules out the substratum answer, on which two objects are two because they have two substrata. We draw three consequences about what does not individuate.

\begin{corollary}[No constituent individuates]\label{c:constituent}
No object is distinct from another in virtue of having a distinct constituent that individuates it.
\end{corollary}

\begin{proof}
An individuating constituent would have to differ from its counterpart in no quality, or it would individuate by its qualities and not as a constituent. Such a constituent is a bare particular, and there are none. Proposition~3.2 of MET-23-a adds that even if one were posited, its distinctness from its counterpart would be primitive, so it would not explain the distinctness of the objects.
\end{proof}

\begin{corollary}[Qualities do not individuate absolutely or relatively]\label{c:qual}
It is not the case that any two objects differ in some qualitative property, or in some qualitative relation in which they stand in one order and not the other.
\end{corollary}

\begin{proof}
This is Corollary~3.3 of MET-28-a: in the world of two qualitatively identical spheres described by \citet{Black1952}, neither form of the principle of the identity of indiscernibles holds.
\end{proof}

\begin{corollary}[Real definitions do not individuate]\label{c:def}
Two objects may have real definitions alike in everything except the object defined.
\end{corollary}

\begin{proof}
A real definition states an object's kind and, where it has one, its origin (MET-13-a, Section~5). Black's spheres are of one kind and may have no origin that distinguishes them, since the world contains nothing else. MET-13-a accordingly denies that real definitions individuate.
\end{proof}

\section{What remains}

These corollaries leave one answer. Two objects are two because they are, and their distinctness is primitive: not grounded in a constituent, a quality or a definition. This is the position Proposition~3.2 of MET-23-a ascribes to the substance view, and the present result is that nothing else is available once substrata are rejected.

Primitive distinctness is compatible with its being witnessed. Black's spheres stand in an irreflexive qualitative relation, being two miles apart, which holds between them and between neither and itself (Proposition~4.1 of MET-28-a). The spheres are weakly discernible, and the bundle theory that MET-29 allows to survive relies on exactly this. Whether the relation grounds the distinctness or presupposes it is disputed, and MET-28-a takes no stand. Our consequence holds either way: if weak discernibility grounds distinctness, what individuates is a relation among the objects themselves, and if it presupposes it, distinctness is primitive. In neither case is a substratum involved.

\begin{corollary}[Haecceities are not substrata]\label{c:haec}
The property of being identical to $a$ exists for every object $a$, but it does not individuate $a$ in the sense in which a substratum was supposed to.
\end{corollary}

\begin{proof}
The property is constructed from identity and $a$, as existence is constructed from identity by quantification (MET-24-a); it presupposes $a$ and its distinctness from everything else, and so cannot be what makes $a$ distinct. A haecceity in the sense of \citet{Adams1979}, a primitive thisness, is the same distinctness stated as a property.
\end{proof}

The last consequence concerns time. A substratum was also supposed to be what persists through change. MET-23-a replies that what persists is the object, and its persistence conditions are fixed by its kind. Diachronic individuation is therefore kind-relative: the ship persists through the replacement of its planks because its kind's persistence conditions are those of a maintained artefact (MET-07), and an organism persists through metabolism because its kind's conditions are biological. No constituent persists beneath these changes and carries the object's identity, since there is none.

The resulting picture is that of a substance ontology without substrata \citep{Lowe2006}: objects bear their properties directly, their distinctness is primitive or witnessed by relations among them, and their identity over time is fixed by what they are.

\begin{thebibliography}{99}\small
\bibitem[Adams(1979)]{Adams1979} Adams, R.~M. (1979). Primitive thisness and primitive identity. \emph{The Journal of Philosophy}, 76(1), 5--26.
\bibitem[Black(1952)]{Black1952} Black, M. (1952). The identity of indiscernibles. \emph{Mind}, 61(242), 153--164.
\bibitem[Lowe(2006)]{Lowe2006} Lowe, E.~J. (2006). \emph{The Four-Category Ontology}. Oxford University Press.
\end{thebibliography}

\end{document}
